\documentclass[11pt]{amsart}

\usepackage{amssymb}
\usepackage{mathtools}
\usepackage{booktabs}
\usepackage{array}
\usepackage[dvipsnames]{xcolor}
\usepackage{listings}
\usepackage{iftex}
\ifXeTeX
  % Under XeTeX the listings `literate` table below is bypassed for non-ASCII
  % input, so the monospaced font itself must contain the Lean symbols.
  \usepackage[no-math]{fontspec}
  \setmonofont{DejaVuSansMono}[Extension=.ttf, UprightFont=*, BoldFont=*-Bold,
    ItalicFont=*-Oblique, BoldItalicFont=*-BoldOblique, Scale=MatchLowercase]
\fi
\usepackage{microtype}
\usepackage[margin=1.1in]{geometry}
\usepackage[hyphens]{url}
\usepackage[
  pdftitle={A formally verified hexadecimal digit extraction for pi with an explicit guard for every position},
  pdfauthor={Brian Bernard Taylor},
  pdfsubject={11Y60, 11J82, 68V20},
  pdfkeywords={pi, BBP formula, digit extraction, irrationality measure, Lean 4, Mathlib, formal verification},
  pdfcreator={LaTeX (Tectonic)},
  colorlinks=true, linkcolor=MidnightBlue, citecolor=OliveGreen, urlcolor=MidnightBlue,
  breaklinks=true
]{hyperref}

%------------------------------------------------------------------
% Theorem environments
%------------------------------------------------------------------
\newtheorem{theorem}{Theorem}[section]
\newtheorem{lemma}[theorem]{Lemma}
\newtheorem{proposition}[theorem]{Proposition}
\newtheorem{corollary}[theorem]{Corollary}
\theoremstyle{definition}
\newtheorem{definition}[theorem]{Definition}
\theoremstyle{remark}
\newtheorem{remark}[theorem]{Remark}

%------------------------------------------------------------------
% Notation
%------------------------------------------------------------------
\newcommand{\N}{\mathbb{N}}
\newcommand{\Z}{\mathbb{Z}}
\newcommand{\R}{\mathbb{R}}
\newcommand{\dn}{d_n}
\newcommand{\Sng}{S_{n,g}}
\newcommand{\Eng}{E_{n,g}}
\newcommand{\lo}{\mathrm{lo}}
\newcommand{\hi}{\mathrm{hi}}
\newcommand{\monus}{\mathbin{\dot{-}}}
\newcommand{\dist}{\delta}
\newcommand{\need}{\nu}
\newcommand{\run}{\rho}
\newcommand{\gap}{\gamma}
\DeclarePairedDelimiter{\floor}{\lfloor}{\rfloor}
\DeclarePairedDelimiter{\abs}{\lvert}{\rvert}
\makeatletter
\g@addto@macro\UrlBreaks{\do\_\do\.}
\makeatother
\DeclareUrlCommand\lean{\urlstyle{tt}}
\setlength{\emergencystretch}{3em}
\newcommand{\Bconst}{1\,795\,710}
\newcommand{\Bsimple}{1\,795\,711}

%------------------------------------------------------------------
% Lean 4 listings
%------------------------------------------------------------------
\definecolor{leankw}{RGB}{0,62,138}
\definecolor{leancomment}{RGB}{110,110,110}
\definecolor{leanbg}{RGB}{248,248,246}
\lstdefinelanguage{Lean4}{
  morekeywords={theorem,lemma,def,noncomputable,abbrev,namespace,end,open,import,
    by,fun,have,show,with,match,if,then,else,let,in,where,instance,structure,Prop,Type},
  sensitive=true,
  morecomment=[l]{--},
  morecomment=[s]{/-}{-/},
  morestring=[b]",
}
\lstset{
  language=Lean4,
  basicstyle=\ttfamily\small,
  keywordstyle=\color{leankw}\bfseries,
  commentstyle=\color{leancomment}\itshape,
  backgroundcolor=\color{leanbg},
  frame=single, rulecolor=\color{black!15}, framesep=4pt,
  xleftmargin=6pt, xrightmargin=6pt,
  columns=fullflexible, keepspaces=true,
  breaklines=true, breakatwhitespace=true,
  showstringspaces=false, upquote=true,
  aboveskip=8pt, belowskip=8pt,
  extendedchars=true,
  literate=
    {ℕ}{{$\mathbb{N}$}}1 {ℤ}{{$\mathbb{Z}$}}1 {ℝ}{{$\mathbb{R}$}}1
    {π}{{$\pi$}}1 {⌊}{{$\lfloor$}}1 {⌋}{{$\rfloor$}}1
    {≤}{{$\le$}}1 {≥}{{$\ge$}}1 {≠}{{$\ne$}}1 {∧}{{$\wedge$}}1 {∨}{{$\vee$}}1
    {¬}{{$\neg$}}1 {→}{{$\to$}}1 {↔}{{$\leftrightarrow$}}1 {∀}{{$\forall$}}1
    {∃}{{$\exists$}}1 {·}{{$\cdot$}}1 {⟨}{{$\langle$}}1 {⟩}{{$\rangle$}}1
    {₁}{{$_1$}}1 {₆}{{$_6$}}1 {∸}{{$\dot{-}$}}1 {∑}{{$\sum$}}1 {∈}{{$\in$}}1
}

%------------------------------------------------------------------
\begin{document}

\title[Verified hexadecimal digits of $\pi$ for every position]
  {A formally verified hexadecimal digit extraction for $\pi$\\
   with an explicit guard for every position}

\author{Brian Bernard Taylor}

\subjclass[2020]{11Y60, 11J82, 68V20}
\keywords{$\pi$, Bailey--Borwein--Plouffe formula, digit extraction,
  irrationality measure, Lean~4, Mathlib, formal verification}

\date{1 November 2026}

\begin{abstract}
We describe a Lean~4 development, built on Mathlib, about the
Bailey--Borwein--Plouffe (BBP) formula for~$\pi$. It contains an independent
proof that the BBP series equals~$\pi$, an integer-only certificate for the
hexadecimal digit $d_n=\floor{16^n\pi}\bmod 16$, and a soundness theorem: when
the two ends of the certificate's bracket agree, the certified digit is the
true digit. To remove the dependence on that check, we formalize an explicit
irrationality-measure bound at the denominators $16^k$: for every $k\ge1$ and
every integer $p$,
$\abs{16^k\pi-p}\cdot16^{570k}>2^{-7182802}\cdot16^k$.
Our main theorem shows that for every $n\ge0$ the certificate returns $d_n$
at every guard depth $g\ge 569n+\floor{\log_{16}(n+1)}+1+B$ with
$B=\Bconst$; a simplified linear form of the guard is also proved. We
also prove a short chain of exact statements relating the run of $0$s or
$\mathrm{F}$s after position~$n$, the distance from $16^n\pi$ to the nearest
integer, and the least sufficient guard depth. All statements are for every
$n$ or for explicit finite ranges; no asymptotic, probabilistic or normality
claim is made. The proofs use only the axioms \lean{propext},
\lean{Classical.choice} and \lean{Quot.sound}, and no \lean{native_decide}.
\end{abstract}

\maketitle

\section{Introduction}\label{sec:intro}

The BBP formula~\cite{BBP1997}
\begin{equation}\label{eq:bbp}
  \pi=\sum_{k\ge0}\frac{1}{16^k}\left(\frac{4}{8k+1}-\frac{2}{8k+4}-\frac{1}{8k+5}-\frac{1}{8k+6}\right)
\end{equation}
lets one compute the hexadecimal digit of $\pi$ at position $n$ without
computing the earlier digits. Every implementation computes a finite sum to
finite precision, so the algorithm needs a \emph{guard}: enough extra digits
that the rounding error cannot change the digit being extracted. How many extra
digits are enough depends on how close $16^n\pi$ is to an integer, which is a
Diophantine question about $\pi$.

This paper reports a Lean~4~\cite{Lean4} development, built on\nMathlib~\cite{Mathlib2020}, that makes the
guard explicit and proves it sufficient for every position. We use the
following conventions throughout. For $n\in\N$ let
\[
  x_n=16^n\pi,\qquad d_n=\floor{x_n}\bmod 16,\qquad
  \dist_n=\min\bigl(\{x_n\},\,1-\{x_n\}\bigr),
\]
so $d_0=3$, $d_1=2$, $d_2=4$, and $\dist_n$ is the distance from $x_n$ to the
nearest integer. The certificate at guard depth $g$ uses
$K=n+g+1$ terms, an integer sum $\Sng\in\N$, an error allowance
$\Eng=4K+2$, and returns $\floor{(\Sng\monus\Eng)/16^g}\bmod16$
(Section~\ref{sec:cert}).

\subsection*{Contributions}
\begin{enumerate}
\item An independent Lean proof of~\eqref{eq:bbp}
  (Theorem~\ref{thm:bbp}). Earlier formal proofs exist in Coq and in
  Lean; see Section~\ref{sec:related}.
\item An integer-only digit certificate with a soundness theorem
  (Theorem~\ref{thm:sound}) and a distance criterion: if
  $2\Eng<16^g\dist_n$ the bracket collapses (Lemma~\ref{lem:dist}).
\item A formal irrationality-measure bound at the denominators $16^k$ with
  exponent $570$ and an explicit constant (Theorem~\ref{thm:measure}).
\item The main theorem, with two co-equal guards: for every $n$ the
  certificate is proved to return $d_n$ at every depth
  $g\ge569n+\floor{\log_{16}(n+1)}+1+B$, $B=\Bconst$ (sharp form), and at
  every depth $g\ge570n+\Bsimple$ (simplified form)
  (Theorem~\ref{thm:main}).
\item Exact relations between the $0$/$\mathrm{F}$ run after position $n$,
  the distance $\dist_n$ and the least sufficient guard depth, and the
  kernel-checked finite statement $\run(77)=0$, $\gap(78)=49$
  (Section~\ref{sec:structure}).
\end{enumerate}

We make no claim about normality or about the statistics of the digits of
$\pi$. Statements of the form ``for every~$n$'' are theorems; statements about
finite ranges of digits are labeled either as kernel-checked theorems or as
computed numerics, and the latter are not used in any proof.

\subsection*{Status of the capped certificate}
The development also contains an earlier capped certificate, which searches
only a fixed set of guard depths; its checked ranges, the unchecked range and
an explicit window where it is proved incorrect are described in
Appendix~\ref{app:future}. It is not part of the main theorem.

\section{Related work}\label{sec:related}

\subsection*{Formal proofs about the BBP formula}
Bertot, Rideau and Th\'ery~\cite{BertotRideauThery2018} formalized in Coq a
proof of the BBP formula and a digit-extraction routine, and used it to
compute the billionth hexadecimal digit of $\pi$ inside Coq; the code is
available as~\cite{TheryPlouffe}. Their correctness statement for the digit
routine is a partial-correctness statement: when the routine returns a digit,
that digit is correct. As the authors note, a routine that never returns a
digit would also satisfy it. They also note that the HOL~Light library contains
a formalization of the BBP formula. Our Theorem~\ref{thm:sound} plays the same
role as their partial-correctness statement; the difference is
Theorem~\ref{thm:main}, which supplies, for every $n$, a guard depth at which
the check is proved to pass.

In Lean, Braams's lean-DLMF project~\cite{BraamsDLMF} contains the theorem
\lean{hasSum_bbp_pi}, the BBP formula in Mathlib's \lean{HasSum} form. That
file was first committed publicly on 16~October 2026. Our proof of
Theorem~\ref{thm:bbp} was developed independently and uses a different
route (an integral representation, see Section~\ref{sec:bbp}). We do not claim
priority for the identity.

\subsection*{\texorpdfstring{Irrationality measures of $\pi$}{Irrationality measures of pi}}
Mahler~\cite{Mahler1953} proved the first explicit bound,
$\abs{\pi-p/q}>q^{-42}$ for all integers $p$ and $q\ge2$. Mignotte~\cite{Mignotte1974}
lowered the exponent to about $20$, and the best published bound we know of in a refereed
journal is $7.103205334137\ldots$ by Zeilberger and Zudilin~\cite{ZeilbergerZudilin2020}.
Our formal exponent $570$ is far from these values; it comes from a
deliberately simple instance of Mahler's method in which every constant is
made explicit and checked in Lean (Section~\ref{sec:measure}). What the
application needs is not a small exponent but an explicit constant valid for
every $k\ge1$.

\subsection*{Other formal work}
Mathlib~\cite{Mathlib2020} proves that $\pi$ is irrational, and Eberl's
Isabelle formalization~\cite{EberlPiTranscendental} proves that $\pi$ is
transcendental. The Prove2Me platform~\cite{Prove2MePi} hosts a campaign of
Lean formulations of irrationality-measure bounds for $\pi$, from Mahler's
exponent $42$ to the value $2$ claimed in a 2026 preprint; on 31~October 2026
its campaign page listed all of these missions as completed. The goal
statements we inspected there are in the form ``for every $\varepsilon>0$ there
is a threshold $Q$'', which asserts the existence of a threshold without
giving one. A guard depth needs a threshold and a constant that are explicit,
so we do not use those results, and we have not audited them.

\section{The BBP identity}\label{sec:bbp}

\begin{theorem}[\lean{bbpSeries_eq_pi}]\label{thm:bbp}
The series~\eqref{eq:bbp} converges and its sum is $\pi$.
\end{theorem}

\begin{proof}[Outline of the formal proof]
Write $p(x)=4-2x^3-x^4-x^5$. Since
$\int_0^1x^{8k+j-1}\,dx=1/(8k+j)$, the $k$-th summand equals
$\int_0^1(x^8/16)^k\,p(x)\,dx$. Dominated convergence for the geometric
series gives
\[
  \sum_{k\ge0}\frac{1}{16^k}\Bigl(\cdots\Bigr)=\int_0^1\frac{p(x)}{1-x^8/16}\,dx .
\]
The factorizations $16-x^8=-(x^2-2)(x^2+2)(x^2+2x+2)(x^2-2x+2)$ and
$p(x)=-(x-1)(x^2+2)(x^2+2x+2)$ reduce the integrand to
$16(x-1)/\bigl((x^2-2)(x^2-2x+2)\bigr)$, which has the antiderivative
$F(x)=2\log(2-x^2)-2\log(x^2-2x+2)+4\arctan(x-1)$ on $[0,1]$. Then
$F(1)-F(0)=\pi$ by the fundamental theorem of calculus.
\end{proof}

\begin{lstlisting}
theorem bbpSeries_eq_pi : bbpSeries = Real.pi
\end{lstlisting}

\section{The digit certificate and its soundness}\label{sec:cert}

\begin{definition}[certificate]\label{def:cert}
For $k,N\in\N$ let
\[
  t_k(N)=\floor*{\frac{4\cdot16^N}{16^k(8k+1)}}\monus
  \floor*{\frac{2\cdot16^N}{16^k(8k+4)}}\monus
  \floor*{\frac{16^N}{16^k(8k+5)}}\monus
  \floor*{\frac{16^N}{16^k(8k+6)}},
\]
where $\monus$ is truncated subtraction on $\N$. For $n,g\in\N$ put
$K=n+g+1$ and
\[
  \Sng=\sum_{k<K}t_k(n+g)\in\N,\qquad \Eng=4K+2,\qquad
  \lo=\floor*{\frac{\Sng\monus\Eng}{16^g}},\qquad
  \hi=\floor*{\frac{\Sng+\Eng}{16^g}} .
\]
The guard \emph{agrees at depth $g$} if $\lo=\hi$. The certified floor is
$\lo$ and the certified digit is $\lo\bmod16$.
\end{definition}

In Lean these are \lean{bbpSpigotTerm}, \lean{bbpGuardedSum},
\lean{bbpGuardErr}, \lean{bbpGuardedLo}, \lean{bbpGuardedHi},
\lean{BBPGuardAgreesAt}, \lean{bbpCertFloorAt} and \lean{bbpCertDigitAt}.
Everything is computed in $\N$. The lemma \lean{bbp_term_cast_bounds} shows
that the truncated subtractions in $t_k(N)$ never truncate and that
\[
  16^N a_k-1<t_k(N)<16^N a_k+3,
\]
where $a_k$ is the $k$-th summand of~\eqref{eq:bbp}. Summing over $k<K$ and
bounding the tail of the series by $1$ gives
\[
  \lo\le x_n<\hi+1\qquad(\text{\lean{bbpGuarded_bracket}}).
\]

\begin{theorem}[soundness; \lean{BBPGuardAgreesAt.floor_pi},
\lean{BBPGuardAgreesAt.digit_pi}]\label{thm:sound}
For all $n,g\in\N$, if the guard agrees at depth $g$, then
$\floor{x_n}=\lo$ and $\lo\bmod16=d_n$.
\end{theorem}

\begin{proof}
If $\lo=\hi$ the bracket gives $\lo\le x_n<\lo+1$, so $\lo=\floor{x_n}$;
Theorem~\ref{thm:bbp} replaces the series by $\pi$.
\end{proof}

Agreement is decidable, so for any particular $(n,g)$ the hypothesis can be
discharged by evaluation in the kernel. The next lemma gives a sufficient
condition in terms of $\dist_n$.

\begin{lemma}[\lean{guardAgreesAt_of_dist}]\label{lem:dist}
Let $n,g\in\N$. If $2\Eng<16^g\abs{x_n-m}$ for every $m\in\Z$, then the guard
agrees at depth $g$. In particular this holds when $2\Eng<16^g\dist_n$.
\end{lemma}

The factor $2$ accounts for one $\Eng$ from the error of $\Sng$ and one from
the width of the bracket. Since $2\Eng=8n+8g+12$ grows linearly in $g$ while
$16^g$ grows exponentially, Lemma~\ref{lem:dist} reduces the choice of a guard
to a lower bound for $\dist_n$.

\section{An irrationality measure at the denominators \texorpdfstring{$16^k$}{16\^{}k}}\label{sec:measure}

\begin{theorem}[\lean{BBPPi.MahlerSharp7b.pi_pow16_irrationality_measure}]\label{thm:measure}
For every integer $k\ge1$ and every $p\in\Z$,
\[
  \abs{16^k\pi-p}\cdot16^{570k}>2^{-7182802}\cdot16^k .
\]
\end{theorem}

\begin{lstlisting}
theorem BBPPi.MahlerSharp7b.pi_pow16_irrationality_measure :
    ∀ k : ℕ, 1 ≤ k → ∀ p : ℤ,
      |(16 : ℝ) ^ k * Real.pi - p| * (16 : ℝ) ^ (k * 570)
        > (1 / 2 ^ (7182801 + 1)) * (16 : ℝ) ^ k
\end{lstlisting}

Equivalently, with $q=16^k$, $\abs{\pi-p/q}>2^{-7182802}q^{-570}$; so the
statement is an irrationality-measure bound with exponent $570$ and an explicit
constant, restricted to the denominators that the digit application uses.

\begin{proof}[Structure of the formal proof]
The proof follows Mahler's method~\cite{Mahler1953}, in the form used by
Mignotte~\cite{Mignotte1974}, with every constant explicit. We list the
ingredients in the order in which they are combined.
\begin{enumerate}
\item \emph{Chebyshev bound.} From Mathlib's bound for the Chebyshev function
  $\psi$ and the estimate $2\sqrt{x}\log x\le 0.02x+8760$ for $x\ge1$ we
  obtain $\psi(x)\le(\log4+0.02)x+8760$ (\lean{psi_le_lin7}), hence
  $\operatorname{lcm}(1,\dots,n)\le\exp(1.4063\,n+8760)$
  (\lean{lcm_le_exp_lin7}). Here $1.4063$ is a rational upper bound for
  $\log 4+0.02$.
\item \emph{Remainder.} For $n\ge1000$ the auxiliary remainder terms satisfy a
  tail bound of the form $(8.55/(m+1))^J$ (\lean{remainder_sharp7}); this uses
  $\pi e<8.5398$. Specialized to $m=141$ and $n=4k+12600$, it gives the
  smallness hypothesis of the next step for every $h\le141$
  (\lean{remainder_spec7}).
\item \emph{Mahler--Mignotte step.} A nonvanishing determinant argument
  (\lean{pi_lower_bound_sharp}) shows that whenever $\abs{\pi-p/q}\le1$,
  \[
    1\le(2q)^{m}\,\abs{\pi-p/q}\;T,
  \]
  where $T$ collects the size of the auxiliary coefficients.
\item \emph{Size of $T$.} With $m=141$ and $n=4k+12600$, the bound
  $T\le2^{7182660+1716k}$ holds for every $k$ (\lean{T_spec7}).
\item \emph{Exponent arithmetic.} With $q=16^k$ we have
  $(2q)^{141}=2^{141+564k}$, so
  $\abs{\pi-p/q}\ge2^{-(141+7182660)}\cdot2^{-(564+1716)k}$.
  Since $141+7182660=7182801$ and $564+1716=2280=4\cdot570$, this is
  $2^{-7182801}\,16^{-570k}$. Passing to the strict inequality costs one
  more factor of $2$, which gives the constant $2^{-7182802}$. The case
  $\abs{\pi-p/q}>1$ is immediate.
\end{enumerate}
The remaining parameters ($m=141$, the offset $12600$, the constants $0.02$
and $8760$) were chosen numerically to keep the exponent and the constant
moderate; the choice affects only the size of the constants, never
correctness.
\end{proof}

\begin{remark}
The same development contains an earlier bound with exponent $658$ and
constant $2^{-1720001}$ (\lean{MahlerSharp.pi_pow16_measure6}), used in
Remark~\ref{rem:658}.
\end{remark}

\section{The main theorem: an explicit guard for every position}\label{sec:main}

Combining Lemma~\ref{lem:dist} with Theorem~\ref{thm:measure} gives a guard
depth that is sufficient for every $n$. Define
\[
  g(n)=569n+\floor{\log_{16}(n+1)}+1+B,\qquad B=\Bconst .
\]

\begin{theorem}[main theorem; \lean{BBPPi.floor_digit_eq_guard570},\\
  \lean{BBPPi.bbpDigit_eq_of_guard_ge}]\label{thm:main}
\leavevmode
\begin{enumerate}
\item \emph{Sharp form.} For every $n\in\N$ the guard agrees at depth
  \[
    g(n)=569n+\floor{\log_{16}(n+1)}+1+\Bconst ,
  \]
  so $\floor{16^n\pi}=\lo(n,g(n))$ and $\lo(n,g(n))\bmod16=d_n$; the same
  holds at every depth $g\ge g(n)$ (\lean{BBPPi.guardAgreesAt_of_guard570_le}).
\item \emph{Simplified form.} For all $n,g\in\N$ with $g\ge570n+\Bsimple$,
  the certificate at depth $g$ returns $\floor{16^n\pi}$ and $d_n$.
\end{enumerate}
\end{theorem}

Both parts are proved in Lean as separate theorems. Part~(1) gives the
smaller depth; part~(2) has a simpler formula. Since
$\floor{\log_{16}(n+1)}+1\le n+1$, we have $g(n)\le570n+\Bsimple$.

\begin{lstlisting}
theorem BBPPi.floor_digit_eq_guard570 (n : ℕ) :
    ⌊(16 : ℝ) ^ n * Real.pi⌋ =
        (bbpCertFloorAt n (569 * n + (Nat.log 16 (n + 1) + 1) + 1795710) : ℤ) ∧
      bbpCertDigitAt n (569 * n + (Nat.log 16 (n + 1) + 1) + 1795710)
        = bbpHexDigitReal n

theorem BBPPi.bbpDigit_eq_of_guard_ge (n g : ℕ) (hg : 570 * n + 1795711 ≤ g) :
    ⌊(16 : ℝ) ^ n * Real.pi⌋ = (bbpCertFloorAt n g : ℤ) ∧
      bbpCertDigitAt n g = bbpHexDigitReal n
\end{lstlisting}

\begin{proof}[Outline]
Write $g=569n+L+j$ with $L=\floor{\log_{16}(n+1)}+1$, so that
$n+1\le16^L$. The generic lemma \lean{guard_agrees_of_const} shows that
$\mathrm{BBPGuardAgreesAt}(n,g)$ holds for every $n$ as soon as $j\ge1$ and
$(8\cdot570+20)(j+1)\le2^{-7182802}\,16^{j}$. For $n\ge1$ it combines
Theorem~\ref{thm:measure} (with $k=n$), which gives
$\dist_n>2^{-7182802}\,16^{-569n}$, with Lemma~\ref{lem:dist}; for $n=0$ it
uses a direct lower bound for the distance from $\pi$ to the nearest
integer. The constant is computed exactly in Lean:
$\ell=\floor{\log_{16}\lceil4580\cdot2^{7182802}\rceil}+1=1795704$
(\lean{guardLogL_570b}), since $4580\cdot2^{7182802}=18320\cdot16^{1795700}$,
and $j=\ell+\floor{\log_{16}(2\ell)}+1=1795704+5+1=\Bconst$
(\lean{guardLogConst_570b}). Part~(2) follows from part~(1), monotonicity
in $g$ and $\floor{\log_{16}(n+1)}\le n$.
\end{proof}

\begin{remark}[the exponent-$658$ guard]\label{rem:658}
The earlier bound gives the guard
$657n+\floor{\log_{16}(n+1)}+1+430009$ (\lean{BBPPi.floor_digit_eq_guard658}), also
proved for every $n$. It exceeds $g(n)$ by $88n-1365701$, so it is the smaller
of the two for $n\le15519$ and the larger for $n\ge15520$. Both theorems
hold for every $n$, so one can take the minimum.
\end{remark}

\subsection*{How far the proved guard is from what is needed}
Let $\need(n)$ be the least $g$ with $2\Eng<16^g\dist_n$ (in Lean,
\lean{piNeed}; it exists because $\pi$ is irrational). Table~\ref{tab:guard}
compares the proved guards with $\need(n)$ at a few positions. The column
$\need(n)$ is \emph{computed numerics}, obtained with \texttt{mpmath} at
precision $4n+400$ bits; it is not a Lean theorem and no result in this paper
depends on it. The other columns are values of the formulas in
Theorem~\ref{thm:main} and Remark~\ref{rem:658}.

\begin{table}[ht]
\centering
\caption{Proved guard depths and the least sufficient depth $\need(n)$.
The last column is computed numerics, not a Lean theorem.}
\label{tab:guard}
\begin{tabular}{@{}r r r r r@{}}
\toprule
$n$ & $g(n)$ (Thm.~\ref{thm:main}(1)) & $570n+\Bsimple$ (Thm.~\ref{thm:main}(2)) & exponent $658$ & $\need(n)$ (numerics)\\
\midrule
$0$       & 1\,795\,711   & 1\,795\,711   & 430\,010      & 2\\
$1$       & 1\,796\,280   & 1\,796\,281   & 430\,667      & 2\\
$10$      & 1\,801\,401   & 1\,801\,411   & 436\,580      & 3\\
$100$     & 1\,852\,612   & 1\,852\,711   & 495\,711      & 4\\
$1000$    & 2\,364\,713   & 2\,365\,711   & 1\,087\,012   & 4\\
$10^4$    & 7\,485\,714   & 7\,495\,711   & 7\,000\,013   & 5\\
$10^5$    & 58\,695\,715  & 58\,795\,711  & 66\,130\,014  & 6\\
$10^6$    & 570\,795\,715 & 571\,795\,711 & 657\,430\,014 & 7\\
\bottomrule
\end{tabular}
\end{table}

The proved guards are very far from $\need(n)$. They are meant as a
correctness guarantee for every position, not as a practical precision
choice: in practice one runs the decidable check of Theorem~\ref{thm:sound}
at a small depth, and Theorem~\ref{thm:main} guarantees that a depth at which
the check passes exists and is at most $g(n)$.

\section{Runs, distance and need}\label{sec:structure}

This section relates three quantities attached to a position $n$. All
results are theorems for every $n$, except where an explicit range is stated.

\begin{definition}
Let $\run(n)$ be the length of the maximal block of equal digits starting at
position $n+1$ that consists only of $0$s or only of $\mathrm{F}$s
(\lean{piRun}; $\run(n)=0$ if $d_{n+1}\notin\{0,\mathrm{F}\}$). For
$s\in\N$, if some digit $d_{s+k}$ with $k\ge0$ equals $1$, let $\gap(s)$ be the
least such $k$ (\lean{piGapSeq}).
\end{definition}

The block in $\run(n)$ is finite because $\pi$ is irrational, so the Lean
definition uses \lean{Nat.find} together with Mathlib's \lean{irrational_pi}.
The Lean definition of $\gap(s)$ returns $0$ when no later $1$ exists. Every
statement below that involves $\gap$ carries the hypothesis that some later
digit equals $1$; we do not prove that this hypothesis holds for every $s$.

\begin{theorem}[\lean{piDist_run}]\label{thm:distrun}
For every $n$, \ $16^{-(\run(n)+1)}<\dist_n<16^{-\run(n)}$.
\end{theorem}

\begin{theorem}[\lean{piNeed_bounds}]\label{thm:needbounds}
For every $n$, with $c_n=\lceil\log_{16}(8n+8\need(n)+12)\rceil$,
\[
  \run(n)+c_n\le\need(n)\le\run(n)+c_n+1 .
\]
\end{theorem}

In Lean $c_n$ is \texttt{Nat.clog~16 (2~*~bbpGuardErr~n (piNeed~n))}, and
$2E_{n,\need(n)}=8n+8\need(n)+12$. Both inequalities in
Theorems~\ref{thm:distrun} and~\ref{thm:needbounds} are sharp in the sense
that equality cannot occur in the strict ones: $16^k\dist_n\ne1$ because
$\pi$ is irrational.

\begin{theorem}[\lean{forced_chain}]\label{thm:chain}
For every $n$:
\begin{enumerate}
\item the digits $d_{n+1},\dots,d_{n+\run(n)}$ are not equal to $1$;
\item if some digit $d_{n+1+k}$ $(k\ge0)$ equals $1$, then
  $\run(n)\le\gap(n+1)$;
\item $16^{-(\run(n)+1)}<\dist_n<16^{-\run(n)}$;
\item $\run(n)+c_n\le\need(n)\le\run(n)+c_n+1$;
\item the guard agrees at every depth $g\ge\need(n)$.
\end{enumerate}
\end{theorem}

\begin{corollary}[\lean{long_run_forces}]
If $L\le\run(n)$, then the digits $d_{n+1},\dots,d_{n+L}$ are not $1$; if
moreover some later digit equals $1$, then $\gap(n+1)\ge L$; and
$\dist_n<16^{-L}$ and $\need(n)\ge L+c_n$.
\end{corollary}

So a long run forces a long gap before the next $1$, a small distance and a
large need. The following explicit finite statement is kernel-checked.

\begin{proposition}[\lean{long_gap_short_run}]\label{prop:converse}
$\run(77)=0$ and $\gap(78)=49$.
\end{proposition}

\begin{proof}
The lemma \lean{cert_78_127} checks, by \mbox{\texttt{decide +kernel}}, that the
certificate's guard agrees at positions $78,\dots,127$, that $d_{78}=3$, that
none of $d_{78},\dots,d_{126}$ equals $1$, and that $d_{127}=1$.
Theorem~\ref{thm:sound} turns these checks into statements about the true
digits.
\end{proof}

By Proposition~\ref{prop:converse}, position $77$ has a gap of $49$ digits
before the next $1$ and no run of $0$s or $\mathrm{F}$s, so, by
Theorem~\ref{thm:distrun}, $\dist_{77}>1/16$. Supporting lemmas in the same file include the countdown
identities \lean{piRun_succ}, \lean{piRun_add} and \lean{piGapSeq_countdown},
the bound \lean{piRun_le_inter_one_gap}, and \lean{piRun_pos_cases}, which
records that a nonempty $0$-run means $x_n$ lies just above an integer and an
$\mathrm{F}$-run that it lies just below one.

\section{The Lean development}\label{sec:lean}

\subsection*{Toolchain}
The development uses Lean \lean{leanprover/lean4:v4.29.0-rc8} and Mathlib at
commit \lean{816ff4582a01b10567d1d4ecb10be7c4ff92aab4}, as recorded in
\lean{lake-manifest.json}. The headline theorems of
Sections~\ref{sec:measure} and~\ref{sec:main} are in the namespace
\lean{BBPPi} (their full names are shown); all other proofs are in the
namespace \lean{FinitistProved} unless another namespace is shown. The
earlier names of the renamed theorems (for example
\lean{bbp_digit_all_n_mu570b}) remain as deprecated aliases in
\lean{Thoughts/Deprecated.lean}. The results of
Section~\ref{sec:structure} are in \lean{FinitistProved.PiDigitStructure},
except the definition \lean{piGapSeq}, which is in
\lean{Thoughts/Proved/BBPDigitDefs.lean}.

\subsection*{Axioms and decision procedures}
For each theorem in Table~\ref{tab:names}, \mbox{\texttt{\#print axioms}} reports no
axioms beyond \lean{propext}, \lean{Classical.choice} and \lean{Quot.sound}
(\lean{cert_78_127} uses only \lean{propext}); in particular there is no
\lean{sorryAx}. The file
\lean{Thoughts/Axioms.lean} runs \mbox{\texttt{\#print axioms}} on every theorem
in the table that is in the default build, and the repository's
continuous-integration job fails if any other axiom, any \lean{sorry}, or
any \lean{native_decide} appears. The theorem of Appendix~\ref{app:future}
for $n\le10000$ is audited separately by \lean{Thoughts/AxiomsHeavy.lean}
after the heavy build. Finite
computations are checked with \mbox{\texttt{decide +kernel}}, so they are evaluated by
the kernel; \lean{native_decide}, which would add the compiler to the trusted
base, is not used.

\subsection*{Theorem names}
Table~\ref{tab:names} lists the Lean name and the file for each numbered
result. All of them are on the \texttt{main} branch of the repository
\texttt{metasecdev/math-proofs} (currently private) at the release tag
\texttt{v1.0.0}, which contains only the BBP/$\pi$ development.

\begin{table}[ht]
\centering
\small
\caption{Results and their Lean names in \texttt{metasecdev/math-proofs},
release tag \texttt{v1.0.0}. File paths are relative to
\texttt{Thoughts/}. Theorem~\ref{thm:cap10000} is in the separate Lake
library \texttt{BBPHeavyCerts}, which GitHub CI does not build; it was
checked by a local build on a MacBook Pro at commit \texttt{d8e418e}; the
modules it depends on are unchanged at \texttt{v1.0.0}.}
\label{tab:names}
\footnotesize
\begin{tabular}{@{}>{\raggedright}p{0.12\textwidth} >{\raggedright}p{0.45\textwidth} >{\raggedright\arraybackslash}p{0.37\textwidth}@{}}
\toprule
Result & Lean name & File\\
\midrule
Thm.~\ref{thm:bbp} & \lean{bbpSeries_eq_pi} & \lean{Proved/BBPPiIdentity.lean}\\
Thm.~\ref{thm:sound} & \lean{BBPGuardAgreesAt.floor_pi}, \lean{BBPGuardAgreesAt.digit_pi} & \lean{Proved/BBPExplicitGuard.lean}\\
Lem.~\ref{lem:dist} & \lean{guardAgreesAt_of_dist} & \lean{Research/BBPConditionalAllN.lean}\\
Thm.~\ref{thm:measure} & \lean{BBPPi.MahlerSharp7b.pi_pow16_irrationality_measure} & \lean{Research/MahlerSharp7b.lean}\\
Thm.~\ref{thm:main}(1) & \lean{BBPPi.guardAgreesAt_guard570}, \lean{BBPPi.floor_digit_eq_guard570}, \lean{BBPPi.guardAgreesAt_of_guard570_le} & \lean{Research/BBPMu570b.lean}\\
Thm.~\ref{thm:main}(2) & \lean{BBPPi.bbpDigit_eq_of_guard_ge} & \lean{Research/BBPMu570b.lean}\\
Rem.~\ref{rem:658} & \lean{BBPPi.floor_digit_eq_guard658} & \lean{Research/BBPMu658.lean}\\
Thm.~\ref{thm:distrun} & \lean{piDist_run} & \lean{Research/PiDigitStructure.lean}\\
Thm.~\ref{thm:needbounds} & \lean{piNeed_bounds} & same file\\
Thm.~\ref{thm:chain} & \lean{forced_chain}, \lean{long_run_forces} & same file\\
Prop.~\ref{prop:converse} & \lean{long_gap_short_run}, \lean{cert_78_127} & same file\\
Thm.~\ref{thm:cap10000} & \lean{BBPSpigotDigitCorrect_of_le_10000_unconditional} & \lean{Proved/BBPDigitsUnconditionalExt10000.lean} (local build, \texttt{d8e418e})\\
Thm.~\ref{thm:capwrong} & \lean{not_BBPSpigotDigitCorrect_window'}, \lean{not_BBPSpigotDigitCorrect_window} & \lean{Research/BBPCapBracket.lean}, \lean{Research/BBPCapRefutation.lean}\\
\bottomrule
\end{tabular}
\end{table}

\subsection*{Dependencies}
The import closure of \lean{PiDigitStructure.lean} contains seven modules of
the development (the BBP algebra, series, tail, guarded-certificate and
identity files, a definitions file, and the file itself). The import closures
of the main theorem and of the capped-certificate results contain only
modules of the BBP/$\pi$ development (the certificate, the irrationality
measure and the explicit Chebyshev bounds); the repository contains no
unrelated modules.

\subsection*{Repository}
The repository \texttt{metasecdev/math-proofs} pins Mathlib in the Lake
configuration as well as in \lean{lake-manifest.json}. Its continuous
integration builds the default target, which contains every result in
Table~\ref{tab:names} except Theorem~\ref{thm:cap10000}, and runs the axiom
check described above. The heavy certificates for Theorem~\ref{thm:cap10000}
need more memory than hosted runners provide and are built separately
(\texttt{lake build BBPHeavyCerts}). The version cited here is the
annotated tag \texttt{v1.0.0}; the repository is private for now.

\section*{Acknowledgements}
The author used AI assistants (Grok Bot) as tools for drafting Lean proofs,
running computations and editing this text. All statements were checked by
the Lean kernel or, where marked as numerics, by independent computation;
the author is responsible for the content.

\appendix

\section{Future work: the capped certificate}\label{app:future}

\subsection{Definition}
The development's original certificate does not use an explicit guard. It
tries the depths $4,8,12,\dots,64$ in order and stops at the first depth at
which the guard agrees; if none agrees it uses depth $68$
(\lean{bbpCertGuard}, \lean{bbpGuardSearch}). Its digit at position $n$ is
then compared with $d_n$; the statement that they are equal is
\mbox{\texttt{BBPSpigotDigitCorrect n}}. When the search stops at an agreeing depth,
Theorem~\ref{thm:sound} applies and the digit is correct. The question is
what happens at the positions where no depth up to $64$ agrees.

\subsection{What is proved}

\begin{theorem}[checked by a local build]\label{thm:cap10000}
For every $n\le10000$ the capped certificate returns $d_n$.
\end{theorem}

In Lean this is \lean{BBPSpigotDigitCorrect_of_le_10000_unconditional}, in
the library \texttt{BBPHeavyCerts}. The proof uses a fast natural-number
checker (\lean{BBPFastGuard.lean}), proved to agree with the guard test, and
checks in the kernel, in blocks, that the guard agrees for each $n\le10000$
at the depth the search selects; it then applies Theorem~\ref{thm:sound}.
The bound $10000$ reflects the cost of the check, not a mathematical limit.

This kernel check needs more than $10$\,GB of memory per Lean process, so it
is not part of the GitHub continuous-integration build. It was checked by a
local build of \texttt{BBPHeavyCerts} on a MacBook Pro (Intel Core i9-9980HK,
64\,GB) at commit \texttt{d8e418e} on 31 October 2026: the build had no errors
and no \lean{sorry}, and the axiom audit \lean{Thoughts/AxiomsHeavy.lean}
reports only \lean{propext}, \lean{Classical.choice} and \lean{Quot.sound}
(report \texttt{reports/heavy-certs-2026-10-31.md}). The same argument
for $n\le2000$ (\lean{BBPSpigotDigitCorrect_of_le_2000_unconditional}) is
checked in the default build, which continuous integration runs.

\begin{theorem}\label{thm:capwrong}
For every $n$ with $2^{272}-71\le n$ and $5n+348\le14\cdot2^{272}$, the capped
certificate does not return $d_n$. In particular this holds for every
$n\in[2^{272},2^{273}]$.
\end{theorem}

In Lean these are \lean{not_BBPSpigotDigitCorrect_window'} (general window)
and \lean{not_BBPSpigotDigitCorrect_window} (the interval
$[2^{272},2^{273}]$).

The positions $10001\le n\le2^{272}-72$ are unchecked: no Lean statement
about the capped certificate at these positions is proved.

\subsection{Mechanism}
For $n\ge1$ the lower end of the bracket satisfies, at every depth~$g$,
\[
  16^gx_n-(5K+3)-16^g<\lo\cdot16^g\le16^gx_n-(K+2)
  \qquad(\text{\lean{guardedLo_position}}),
\]
with $K=n+g+1$. Thus $\lo$ is biased downward from $x_n$ by at least
$(K+2)/16^g$. For $n\ge2^{253}$ no depth up to $64$ can agree
(\lean{bbpCertGuard_eq_68}), so the search uses $g=68$. Once
$16^{68}\le K+2$, that is, $n\ge2^{272}-71$, the downward bias is at least one
unit, while the upper condition $5K+3\le14\cdot16^{68}$ keeps $\lo$ within the
next lower integer; together they force $\lo\bmod16\ne d_n$
(\lean{lo_digit_ne_of_window}). Below $2^{272}-71$ the bias is smaller than
one unit, and whether the capped digit equals $d_n$ depends on the actual
digits of $\pi$ after position $n$ (roughly, on a run of about
$68-\log_{16}n$ zeros or $\mathrm{F}$s). These positions are unchecked.

\subsection{Directions}
\begin{enumerate}
\item \emph{Sharper constants.} The exponent $570$ and the constant $B$ come
  from the factor $1.4063$ in the bound for $\operatorname{lcm}(1,\dots,n)$
  and from the other explicit choices in Section~\ref{sec:measure}. Sharper
  explicit Chebyshev bounds would lower both.
\item \emph{Kernel-checked tables of $\need(n)$.} The values in the last
  column of Table~\ref{tab:guard} could be turned into kernel-checked
  statements for explicit ranges, using Theorem~\ref{thm:needbounds} and the
  decidable check.
\item \emph{The heavy certificates.} Run the heavy build for
  Theorem~\ref{thm:cap10000} in automated checking, and extend the kernel-checked range of the
  capped certificate beyond $n=10000$.
\end{enumerate}

\bibliographystyle{amsplain}
\bibliography{refs}

\end{document}
